\section{Where the Information Lives}
\label{sec:marginal}

With the backbone fixed, we measure what exogenous information adds to it. The
measurement is made inside a single model class so that attribution is never
confounded with model capacity: Section~\ref{sec:linvsnonlin} varies the class
with the information set held fixed, and this section does the reverse.

\subsection{Benchmark Debt}
\label{sec:benchmark_debt}

Section~\ref{sec:powerlaw} established that the shared incumbent---OLS on the
base-5 HAR ladder plus calendar, QLIKE $0.13415$---is beaten by the same
estimator on a base-2 ladder, QLIKE $0.13332$, at DM $t = -10.6$. That gap of
$0.00082$ involves no additional information, no change of model class, and no
tuning. It is pure quadrature error in the benchmark.

We state the consequence before reporting any increment, because it applies to
all of them. Every $\Delta$QLIKE in this section and in
Section~\ref{sec:linvsnonlin} is measured against the base-5 incumbent, and
therefore \emph{includes} that $0.00082$ of backbone misspecification as though
it were exogenous or nonlinear content. For the smaller effects this is not a
correction but a reclassification: the market-equal-weighted bucket's apparent
$0.00096$ gain is, against a correctly specified linear backbone, on the order
of the quadrature error itself.

The correct remedy is to re-run the bucket battery on the base-2 backbone so
that the panel, the benchmark and the ladder convention are unified.
\pending{bucket battery on the b2 backbone; this is a correctness prerequisite,
not a robustness extension, and until it lands every increment in
Tables~\ref{tab:buckets} and~\ref{tab:class} should be read as an upper bound
carrying $0.00082$ of benchmark debt}

Two mitigating observations. First, the debt is a constant additive offset
shared by every arm, so \emph{orderings} within a table are unaffected; it is
levels and the significance of the smallest effects that move. Second, the
largest effects in the paper---the joint exogenous gain of $0.00627$ and the
model-class results of Section~\ref{sec:linvsnonlin}---survive the correction
comfortably, while the bottom half of Table~\ref{tab:buckets} does not survive
it cleanly.

\subsection{Every Bucket Helps, and They Heavily Overlap}
\label{sec:buckets}

Table~\ref{tab:buckets} reports the causally tuned ridge arm on each of the
eight exogenous buckets, on the shared panel, against the shared incumbent.
Ridge is the attribution instrument: it is the dense-shrinkage estimator that
Section~\ref{sec:linvsnonlin} shows to be the appropriate linear response to
this regime, and holding it fixed while varying the information set isolates
information from capacity.

\begin{table}[H]
\centering
\caption{Bucket attribution. Causally tuned ridge on each exogenous bucket,
shared panel ($n = 218{,}934$), identical rows for every cell. $\Delta$ is
measured against the shared incumbent (OLS HAR base-5 $+$ calendar,
$0.13415$) and carries the benchmark debt of Section~\ref{sec:benchmark_debt}.
$\Delta_{\text{own}}$ is measured against the same estimator's own no-exogenous
baseline ($0.13445$), which is the correct within-class increment. Negative DM
$t$ favours the bucket arm.}
\label{tab:buckets}
\begin{tabular}{lrrrr}
\toprule
Bucket & QLIKE & $\Delta$ vs.\ incumbent & $\Delta_{\text{own}}$ & DM $t$ \\
\midrule
all\_features (joint) & \textbf{0.12788} & $-0.00627$ & $-0.00657$ & $-19.9$ \\
\midrule
moments          & 0.13085 & $-0.00331$ & $-0.00360$ & $-19.2$ \\
liquidity        & 0.13120 & $-0.00295$ & $-0.00325$ & $-13.4$ \\
implied\_vol     & 0.13252 & $-0.00163$ & $-0.00193$ & $-11.3$ \\
market\_vw       & 0.13291 & $-0.00125$ & $-0.00154$ & $-9.1$ \\
market\_ew       & 0.13319 & $-0.00096$ & $-0.00126$ & $-7.2$ \\
vol\_demand      & 0.13337 & $-0.00078$ & $-0.00108$ & $-4.1$ \\
sentiment        & 0.13347 & $-0.00068$ & $-0.00098$ & $-5.7$ \\
\midrule
(no exogenous)   & 0.13445 & $+0.00030$ & --- & $+5.0$ \\
\bottomrule
\end{tabular}
\end{table}

Two facts organise the table.

\paragraph{Every bucket improves on the backbone, decisively.} All seven
singleton buckets carry DM $t$ between $-4.1$ and $-19.2$. Taken individually
these are overwhelming; Section~\ref{sec:multiplicity} explains why that is
less impressive than it looks.

\paragraph{The buckets overlap almost completely.} Summing the seven
within-class increments gives $0.01364$, against a joint increment of
$0.00657$: a ratio of $2.08$. \emph{The sum of the parts is roughly twice the
whole.} The same computation on the causally tuned LightGBM arm gives $2.24$,
so the redundancy is a property of the information rather than of the linear
estimator.

We note a trap here that the earlier drafts of this analysis fell into. If the
overlap ratio is computed against the \emph{shared incumbent} rather than each
estimator's own baseline, every bucket's increment silently includes that
estimator's model-class premium, and the ratio inflates---to $3.98$ for
LightGBM, against its true $2.24$. Overlap ratios must be computed within
class.

The practical reading is that the eight buckets are eight noisy views of a
small number of underlying states, not eight independent information sources.
This is the same fact as the dense-but-weak finding of
Section~\ref{sec:dense_weak}, observed one level up, and it is what the
low-rank operator structure of Section~\ref{sec:hawkes} would predict.

\subsection{Block Attribution}
\label{sec:fwl}

A Frisch--Waugh--Lovell block decomposition assigns each block its Type-III
unique contribution---its gain net of everything else in the design.

\begin{table}[H]
\centering
\caption{FWL block attribution of the linear base. Each entry is the block's
unique contribution net of all other blocks. \textbf{Panel note:} this
decomposition was computed on the earlier base-5 exogenous cache (full fit
QLIKE $0.12314$) and is \emph{not} on the frozen battery panel; it is reported
here for its structure, not for level comparability with
Table~\ref{tab:buckets}.}
\label{tab:fwl}
\begin{tabular}{lrr}
\toprule
Block & Columns & Unique contribution \\
\midrule
HAR backbone                  & 6   & $-0.02208$ \\
Exogenous                     & 359 & $-0.00449$ \\
Regime (HAR $\times$ \{open, close\}) & --- & $-0.00225$ \\
\texttt{cumrv} (engineered)   & 1   & $-0.00122$ \\
\bottomrule
\end{tabular}
\end{table}

The backbone accounts for 96\% of explainable variance ($R^2 = 0.594$ of
$0.617$ for the full fit). Everything built on top of it is small. The
per-feature comparison is the one the attribution literature would find
striking: the single engineered path feature \texttt{cumrv} contributes
$-0.00122$ on its own, while the 359-column raw exogenous block contributes
$-0.00449$ in total, or $1.25 \times 10^{-5}$ per column---a factor of
\begin{equation}
  \frac{0.00122 / 1}{0.00449 / 359} \;\approx\; 97.5
  \label{eq:per_feature}
\end{equation}
in favour of one mechanism-targeted measurement over piling on raw series.

We show the arithmetic because per-feature normalisation is easy to abuse. Two
caveats belong with the number. First, dividing a block's \emph{unique}
(post-FWL) contribution by its column count penalises internal redundancy
twice: the exogenous block is heavily collinear, so its Type-III contribution
is small precisely because its members substitute for one another, and then
that small number is divided by the large count that caused the substitution.
Second, \texttt{cumrv} was engineered \emph{after} inspecting this panel,
whereas the raw exogenous series were not; the comparison is between a
searched-for feature and an unsearched block. The honest claim is directional:
mechanism-targeted measurement is far more efficient per column than
indiscriminate accumulation of price-derived series. The factor of 97.5 should
not be quoted as an effect size.

\paragraph{Discrepancy note.} An earlier internal memo records this ratio as
``$\sim$27$\times$''. Recomputing from the block contributions above gives
$97.5$. The memo figure appears to be an error and should not be propagated.

\subsection{The Model Confidence Set Over Feature Sets}
\label{sec:mcs}

Earlier drafts asserted that the 90\% model confidence set
\citep{HansenLundeNason2011} over the nine feature sets is the singleton
containing the joint bucket. \textbf{We withdraw that claim.} A model
confidence set requires the arms' per-bar loss series, which were not retained
for the battery, so the assertion was never supported by a computed set;
Section~\ref{sec:mcs_results} reports the confidence sets we \emph{can}
compute, on the campaign tile, where they behave quite differently.
\pending{battery MCS over the nine feature sets once per-bar losses are
persisted}

Two interpretive cautions will survive the eventual computation. First, with
$n = 218{,}934$ bar-level loss differentials an MCS has enormous power, and
collapsing to a singleton is close to the default outcome whenever one arm is
even slightly and systematically ahead; such a result would mean ``the joint
bucket is not statistically tied with any singleton bucket'', not ``the joint
bucket is the only useful feature set''---a distinction
Table~\ref{tab:buckets}'s overlap ratio already makes on economic grounds.
Second, the tile results run the other way: at $n = 2{,}189$ the MCS retains
123 of 161 arms and all 20 of the best. Sample size, not model quality, drives
how discriminating a confidence set looks, and neither extreme should be
reported without the other.

\subsection{Multiplicity}
\label{sec:multiplicity}

Every DM statistic in Table~\ref{tab:buckets} is computed against one shared
comparator, and the full battery from which the table is drawn contains 98
arms, all compared to the same incumbent, with the best reported as the
headline. No multiplicity adjustment is applied anywhere in this paper.

Section~\ref{sec:multiplicity_results} now measures the cost. Across all 97
testable arms, 93 are significant at an unadjusted $p < 0.05$; Holm--Bonferroni
leaves 91 and Benjamini--Hochberg leaves 93. Every bucket effect in
Table~\ref{tab:buckets} survives Holm, including \texttt{vol\_demand} at
$t = -4.1$, whose raw $p = 4.3 \times 10^{-5}$ becomes $4.2 \times 10^{-3}$
after correction. Multiple testing is therefore not what threatens the bottom
of this table---the benchmark debt of Section~\ref{sec:benchmark_debt} is, and
that threat is unaffected by any $p$-value adjustment. Of the two arms in the
whole battery that do lose significance under Holm, one is the marginal
\texttt{reclasso/sentiment} arm and the other a base-5 ladder cap variant;
neither carries a claim in this paper.

\subsection{Where the Attribution Points}
\label{sec:attribution_points}

The attribution results point outward rather than inward. The backbone carries
almost all of the explainable variance; the exogenous block adds a real but
small increment that is heavily redundant internally; and a single feature
built to measure a specific intraday mechanism outperforms the raw block on a
per-column basis by a wide margin. The implied next investment is not a larger
model but mechanism-targeted data---auction imbalance, order flow, and dealer
positioning aligned to the session transitions where
Sections~\ref{sec:clock} and~\ref{sec:saturation} both locate the structure.
